\section{Accuracy shrinking with $\delta$}\label{sec:joint}
We now let $K$ tend to zero with $\delta$. The constants in Section~\ref{sec:fixed} depend on the fixed approximation order, so those bounds do not apply when the order grows. We prove two independent lower bounds with constants uniform in the order. When $K\leq\delta^\beta$ for a fixed $\beta>0$, each matches the upper bound of a contractive construction (Proposition~\ref{prop:integrated-sign}), which determines the order of the query count (Theorem~\ref{thm:joint-matched}). For the intermediate regime, where $K\to0$ but $\log(1/K)=o(\log(1/\delta))$, we give a pinned construction with constants uniform in the order. Throughout this section $\rho>1$ is fixed,
\[
 D=\log(1/\delta),\qquad L=\log(1/K),\qquad
 R_0=\frac{\sqrt\rho+1}{\sqrt\rho-1},\qquad a=\log R_0.
\]
Constants denoted by $C_\rho,c_\rho$ may change between occurrences but never depend on a growing degree. When a fixed constant is raised to such a degree, the power is always displayed.

\subsection{Exterior growth with the exact sine target}
For every integer $n\geq2$ set
\begin{equation}\label{eq:Hn}
 H_n=h_\rho\max_{z\geq z_0}\frac{z-z_0}{T_n(z)}.
\end{equation}
Thus $H_{2r}=G_r$; odd indices are also used below. The proof of Proposition~\ref{prop:threshold-asymptotics} applies without a parity assumption and gives
\begin{equation}\label{eq:H-asymptotic}
 H_n=\frac{2\sqrt\rho}{\mathrm e n}R_0^{-n}(1+O_\rho(n^{-1})).
\end{equation}
More precisely, the unique maximizing point is $z_n=\cosh(a+s_n/n)$ with $s_n=1+O_\rho(n^{-1})$, by the stationary equation
$n(\cosh t-\cosh a)\tanh(nt)=\sinh t$ for $z=\cosh t$.
Consequently $\xi_n:=h_\rho(z_n-z_0)=\Theta_\rho(1/n)$. Also $H_n$ strictly decreases to zero: evaluate the ratio for $H_{n+1}$ at its maximizer and use $T_{n+1}(z)>T_n(z)$ for $z>1$.

\begin{theorem}[Quantitative exterior lower bound]\label{thm:joint-exterior}
There are $\delta_0,K_0>0$, depending only on $\rho$, such that the following holds for $0<\delta<\delta_0$ and $0<K<K_0$. Let $n\geq3$ be the largest odd integer such that $H_n\geq8K$. Every coherent converter that uses $T$ queries and is accurate to $K\delta$ on $[\delta,\rho\delta]$ satisfies
\begin{equation}\label{eq:exterior-query}
 T\geq c_\rho nK^{1/(n+1)}\delta^{-n/(n+1)}.
\end{equation}
Moreover,
\begin{equation}\label{eq:nK}
 n=\frac{L-\log L+O_\rho(1)}a,
 \qquad K^{1/(n+1)}=\Theta_\rho(1).
\end{equation}
The constants and $\delta_0$ do not depend on $K$, so $L$ may grow arbitrarily fast.
\end{theorem}
\begin{proof}
On the scalar family of Lemma~\ref{lem:scalar}, put $F(t)=1-\operatorname{Re}q(t)$. Then $F$ has trigonometric degree at most $T$, satisfies $0\leq F\leq2$ globally, and
\[
 |F(\arcsin x)-x|\leq K\delta\qquad(\delta\leq x\leq\rho\delta).
\]
Define
\[
 b_\delta=\frac{\arcsin(\rho\delta)-\arcsin\delta}{\rho-1},
 \qquad a_\delta=\arcsin\delta-b_\delta.
\]
The affine map $y\mapsto a_\delta+b_\delta y$ sends $[1,\rho]$ onto the angle interval $[\arcsin\delta,\arcsin(\rho\delta)]$. Strict convexity of $\arcsin$ on $(0,1)$ gives $b_\delta>\arcsin\delta>\delta$, $a_\delta<0$, and $b_\delta\leq C_\rho\delta$ for small $\delta$.
Let $P_n$ and $S_n$ be the order-$n$ Taylor polynomials, in $y$, of $F(a_\delta+b_\delta y)/\delta$ and $\sin(a_\delta+b_\delta y)/\delta$, respectively. For $|y|\leq A_0:=\rho+1$, Taylor's theorem and the trigonometric Bernstein inequality give remainders bounded by
\begin{equation}\label{eq:joint-remainders}
 u_T=\frac{2(C_\rho A_0T)^{n+1}\delta^n}{(n+1)!},
 \qquad u_s=\frac{(C_\rho A_0)^{n+1}\delta^n}{(n+1)!}.
\end{equation}
We keep the exact sine target: its linear part would cause an error of order $\delta^2$ after rescaling, too large when $K\ll\delta^2$. Correctness therefore implies
\[
 \norm{P_n-S_n}_{[1,\rho]}\leq K+u_T+u_s.
\]
Decrease $K_0$ so that $0<\xi_n\leq1$. Global nonnegativity yields $P_n(-\xi_n)\geq-u_T$. Also
$\arcsin(\delta \xi_n)\leq \xi_n\arcsin\delta$, since $\arcsin x/x$ increases for positive $x$. Hence
\[
 b_\delta(1+\xi_n)\geq\arcsin\delta+\arcsin(\delta \xi_n),
 \qquad a_\delta-b_\delta \xi_n\leq-\arcsin(\delta \xi_n).
\]
All these angles lie in $(-\pi/2,\pi/2)$ for small $\delta$, uniformly in $n$, so $S_n(-\xi_n)\leq-\xi_n+u_s$.
The exterior Chebyshev inequality, proved in Theorem~\ref{thm:thresholds} and classically stated in \cite[Proposition 2.1]{bos2020}, now gives
\[
 \xi_n-u_T-u_s\leq (K+u_T+u_s)T_n(z_n).
\]
Since $T_n(z_n)\geq1$ and $\xi_n/T_n(z_n)=H_n$,
\begin{equation}\label{eq:exterior-obstruction}
 H_n\leq K+2u_T+2u_s.
\end{equation}
Maximality of $n$ gives $8K>H_{n+2}\geq c_\rho(n+2)^{-1}R_0^{-n-2}$, using \eqref{eq:H-asymptotic}. Consequently
\[
 \frac{u_s}{K}\leq
 C_\rho(n+2)\frac{(C_\rho R_0\delta)^n}{(n+1)!}.
\]
For sufficiently small $\delta_0$ this is at most one simultaneously for every allowed $n$. Equations \eqref{eq:exterior-obstruction} and $H_n\geq8K$ then imply $u_T\geq(5/2)K$. Solving \eqref{eq:joint-remainders} and using $((n+1)!)^{1/(n+1)}\geq(n+1)/\mathrm e$ proves \eqref{eq:exterior-query}. Finally, the two inequalities $H_n\geq8K>H_{n+2}$ and \eqref{eq:H-asymptotic} give \eqref{eq:nK}. The proof uses the scalar oracles of Lemma~\ref{lem:scalar} at all angles $t$ but correctness only for $t\in[\arcsin\delta,\arcsin(\rho\delta)]$; it never assumes that the output depends only on $\sin t$.
\end{proof}

\subsection{A contractive upper bound and matching bounds at high accuracy}
\begin{proposition}[Integrated sign construction]\label{prop:integrated-sign}
For $0<\delta<1$ and $0<\eta<1/4$, there is a real polynomial $q$ of degree
$O(\delta^{-1}\log(1/\eta))$ with $0\leq q\leq1$ on $[-1,1]$ and
\[
 |q(x)-(1-x)|\leq\eta\qquad(\delta\leq x\leq1).
\]
Hence there is a unit-normalized converter with the same query order.
\end{proposition}
\begin{proof}
The bounded sign-approximation lemma \cite[Lemma 25, extended version]{gilyen2019} supplies an odd real polynomial $s$ of degree $O(\delta^{-1}\log(1/\eta))$, bounded by one on $[-1,1]$, with $|s(x)-\operatorname{sign}x|\leq2\eta$ for $|x|\geq\delta$. Set
\[
 p(x)=\int_{-1}^x\frac{1+s(t)}2\,dt,\qquad q(x)=1-p(x).
\]
Oddness gives $p(1)=1$, and $0\leq p'\leq1$ gives $0\leq p,q\leq1$. For $x\geq\delta$,
\[
 q(x)-(1-x)=\frac12\int_x^1(s(t)-1)\,dt,
\]
whose absolute value is at most $\eta$. Lemma~\ref{lem:implementation}, which needs no definite parity and loses no normalization, gives the converter.
\end{proof}

\begin{theorem}[Optimal query count at high accuracy]\label{thm:joint-matched}
For each fixed $\beta>0$, uniformly for $0<K\leq\delta^\beta$ and sufficiently small $\delta$ satisfying \eqref{eq:bands},
\begin{equation}\label{eq:matched-shift}
 Q(\delta,K\delta;\rho,c)
 =\Theta_{\rho,\beta}\!\left(\delta^{-1}\log(1/K)\right).
\end{equation}
Let $Q_{\rm comp}(\delta,\eta)$ instead denote the worst-case query count for unit-normalized conversion to $I-H$ with operator-norm error $\eta$, for all $H$ with $\delta I\preceq H\preceq I$. For every fixed $\beta>0$, uniformly for $0<\eta\leq\delta^{1+\beta}$,
\begin{equation}\label{eq:matched-complement}
 Q_{\rm comp}(\delta,\eta)
 =\Theta_\beta\!\left(\delta^{-1}\log(1/\eta)\right).
\end{equation}
\end{theorem}
\begin{proof}
Equation \eqref{eq:exterior-query}, with \eqref{eq:nK}, is at least
$c_\rho\delta^{-1}L\exp(-C_\rho D/L)$. If $L\geq\beta D$, the exponential factor is bounded below by a positive constant depending only on $\rho,\beta$. Proposition~\ref{prop:integrated-sign} gives the upper bound $O(\delta^{-1}(D+L))$, which is $O_\beta(\delta^{-1}L)$. For $Q_{\rm comp}$, the promise includes spectra in $[\delta,2\delta]$, so the lower bound applies with $\rho=2$ and $K=\eta/\delta$. Writing $E=\log(1/\eta)$ gives $L=E-D\geq\beta D$, hence $L\geq\beta E/(1+\beta)$. The upper construction is already accurate on $[\delta,1]$.
\end{proof}
These bounds hold for every dimension and sparsity whenever the promise includes the scalar instances $H=x$ for all $x\in[\delta,\rho\delta]$ (for $Q_{\rm comp}$, all $x\in[\delta,2\delta]$). They need not hold for known finite spectra (Section~\ref{sec:exact}) or when the normalized complement is directly accessible. Standard uniform singular value amplification \cite[Theorem 30, extended version]{gilyen2019} also gives the upper order, after first encoding $(I-H)/2$. The new part is the lower bound, which relies on correctness for a whole interval of eigenvalues and on unit normalization.

\subsection{An independent lower bound via Fej\'er--Riesz factorization}
The second lower bound uses Fej\'er--Riesz factorization instead of exterior extrapolation of a Taylor polynomial, and does not use Theorem~\ref{thm:joint-exterior}.

\begin{lemma}[Fej\'er--Riesz factor and its Taylor remainder]\label{lem:FR-remainder}
For a $T$-query converter put $F(t)=1-\operatorname{Re}q(t)$. There is a polynomial $A$ of degree at most $T$ such that
\[
 F(t)=|A(e^{it})|^2,\qquad \norm{A}_{|z|=1}\leq\sqrt2.
\]
Let $b(x)=A(e^{i\arcsin x})$, analytic for $|x|<1$, and let $B_r$ be its Taylor polynomial of degree $r$. Set $s=r+1$. For $\rho\delta<R\leq1/4$,
\begin{equation}\label{eq:FR-general-remainder}
 e_r:=\sup_{1\leq y\leq\rho}|b(\delta y)-B_r(\delta y)|
 \leq \frac{\sqrt2 e^{2TR}}{1-\rho\delta/R}
             (\rho\delta/R)^s.
\end{equation}
If the converter has error $K\delta$ on the low band, then, with $u=e_r/\sqrt\delta$,
\begin{equation}\label{eq:FR-obstruction}
 G_r\leq K+2\sqrt{\rho+K}\,u+u^2.
\end{equation}
\end{lemma}
\begin{proof}
The first assertion is the classical Fej\'er--Riesz theorem; its root-factorization proof is also given in \cite[proof of Theorem 4]{motlagh2024}. To recall why it applies, multiply the nonnegative Laurent polynomial $F$ of actual degree $d$ by $z^d$. Its roots away from the unit circle occur in conjugate-reciprocal pairs, and its roots on the circle have even multiplicity, because a real analytic nonnegative function has zeros of even order. Selecting one root from each reciprocal pair and half of each circle multiplicity gives $A$; positivity fixes a nonnegative scalar factor. The case $F\equiv0$ is immediate. The bound on its norm on the unit circle follows from $F\leq2$.
The maximum principle applied to $A$ inside the unit disk and to $z^{-T}A(z)$ outside it gives
\[
 |A(z)|\leq\sqrt2\max(1,|z|^T).
\]
The power series for $\arcsin$ has nonnegative coefficients and yields $|\arcsin z|\leq\arcsin R\leq2R$ for $|z|=R\leq1/4$. Thus $|b(z)|\leq\sqrt2e^{2TR}$. Cauchy's coefficient estimate and summation of a geometric tail prove \eqref{eq:FR-general-remainder}.
For real $y$, the polynomial
\[
 P(y)=\delta^{-1}B_r(\delta y)\overline{B_r(\delta y)}
\]
(where conjugation acts on coefficients) is real, globally nonnegative, and of degree at most $2r$. On the low band, $|b(\delta y)|\leq\sqrt{(\rho+K)\delta}$ and $||b(\delta y)|^2/\delta-y|\leq K$. The inequality $||B|^2-|b|^2|\leq2|b||B-b|+|B-b|^2$ and the definition of $G_r$ prove \eqref{eq:FR-obstruction}.
\end{proof}

We call $G_r-K$ the margin of $K$ below the threshold $G_r$.
\begin{theorem}[Uniform finite-index margin bound]\label{thm:FR-margin}
There are $c_\rho,C_\rho>0$ such that, for $0<\delta\leq1/(2\rho)$, $r\geq1$, $K<G_r$, and $r\leq c_\rho\delta^{-1/2}$, every converter accurate to $K\delta$ on the low band satisfies
\begin{equation}\label{eq:FR-margin-bound}
 T\geq C_\rho(r+1)(G_r-K)^{1/(r+1)}
                  \delta^{-1+1/(2r+2)}.
\end{equation}
\end{theorem}
\begin{proof}
First, $G_r\leq G_1<G_0=(\rho-1)/2$. The values of $|b|$ at $x=\delta$ and $x=\rho\delta$, with $b$ as in Lemma~\ref{lem:FR-remainder}, differ by at least
\[
 \bigl(\sqrt{\rho-K}-\sqrt{1+K}\bigr)\sqrt\delta
 \geq d_\rho\sqrt\delta,
\]
with $d_\rho>0$ since $K<G_1$. Bernstein's inequality for the analytic trigonometric polynomial $A(e^{it})$ bounds its derivative by $\sqrt2T$. The two endpoint angles differ by at most $C_\rho\delta$, so $T\geq d'_\rho\delta^{-1/2}$. Decreasing $c_\rho$ ensures $s=r+1\leq T$.
If $T\rho\delta\geq s/8$, the resulting $T\geq s/(8\rho\delta)$ is stronger than \eqref{eq:FR-margin-bound}, because $(G_r-K)^{1/s}\delta^{1/(2s)}$ is bounded above by a $\rho$-dependent constant. Otherwise use $R=s/(4T)\leq1/4$, so $\rho\delta/R<1/2$ in \eqref{eq:FR-general-remainder}. It gives
\begin{equation}\label{eq:FR-adaptive-radius}
 e_r\leq2\sqrt2\left(4\sqrt{\mathrm e}\,T\rho\delta/s\right)^s.
\end{equation}
By \eqref{eq:FR-obstruction},
\[
 u\geq\sqrt{\rho+G_r}-\sqrt{\rho+K}
 =\frac{G_r-K}{\sqrt{\rho+G_r}+\sqrt{\rho+K}}
 \geq c_\rho(G_r-K).
\]
Combine this with \eqref{eq:FR-adaptive-radius} and take an $s$th root. Constants remain uniform because $s\geq2$.
\end{proof}

\begin{theorem}[Lower bound for every admissible index $r$]\label{thm:FR-all}
For sufficiently small $\delta>0$, depending only on $\rho$, let a converter use $T$ queries and be accurate to $K\delta$ on $[\delta,\rho\delta]$. Then every integer $r\geq1$ satisfying
\begin{equation}\label{eq:FR-halfthreshold}
 K\leq\frac{\sqrt\rho}{4\mathrm e r}R_0^{-2r}
\end{equation}
gives the lower bound
\begin{equation}\label{eq:FR-all-bound}
 T\geq c_\rho r\delta^{-1+1/(2r+2)}.
\end{equation}
There is no upper restriction on $r$ or on $\log(1/K)$.
\end{theorem}
\begin{proof}
Evaluating \eqref{eq:G-formula} at $z=\cosh(a+1/(2r))$, using convexity of $\cosh$ and $\cosh(2ra+1)\leq e^{2ra+1}$, gives the explicit estimate
\begin{equation}\label{eq:G-explicit-lower}
 G_r\geq \frac{\sqrt\rho}{2\mathrm e r}R_0^{-2r}.
\end{equation}
Thus \eqref{eq:FR-halfthreshold} implies $G_r-K\geq G_r/2$. Inequality \eqref{eq:FR-obstruction} of Lemma~\ref{lem:FR-remainder} consequently requires
\begin{equation}\label{eq:FR-required-error}
 e_r\geq c_\rho\sqrt\delta\,r^{-1}R_0^{-2r}.
\end{equation}
If $T<s=r+1$, choose the fixed radius $R=1/4$. For $\rho\delta\leq1/8$, \eqref{eq:FR-general-remainder} gives
\[
 e_r\leq2\sqrt2 e^{s/2}(4\rho\delta)^s.
\]
The ratio of this upper bound to the right side of \eqref{eq:FR-required-error} is at most
\[
 C_\rho r\delta^{r+1/2}
             (4\rho\sqrt{\mathrm e}\,R_0^2)^r.
\]
It is less than one for every $r\geq1$ simultaneously when $\delta$ is sufficiently small: choose $C'_\rho\delta<1/2$ and use $\sup_{r\geq1}r2^{-r}\leq1$. Hence $T\geq s$. The argument of Theorem~\ref{thm:FR-margin} with radius $R=s/(4T)$ now applies. There, the restriction $r\leq c_\rho\delta^{-1/2}$ was used only to ensure $T\geq s$. Taking an $s$th root of \eqref{eq:G-explicit-lower} costs only a fixed factor, since $r^{-1/(r+1)}$ is bounded below and $R_0^{-2r/(r+1)}\geq R_0^{-2}$. This proves \eqref{eq:FR-all-bound}.
\end{proof}
In particular, for $0<K\leq\sqrt\rho R_0^{-2}/(4\mathrm e)$, the following maximum is over a finite nonempty index set and gives the explicit lower bound
\begin{equation}\label{eq:FR-max}
 T\geq c_\rho\max_{\substack{r\geq1:\
 K\leq\sqrt\rho R_0^{-2r}/(4\mathrm e r)}}
 r\delta^{-1+1/(2r+2)}.
\end{equation}
For large $L$, the largest admissible integer is
$r=(L-\log L+O_\rho(1))/(2a)$, obtained by taking logarithms in \eqref{eq:FR-halfthreshold}. Alternatively, choosing $r=\lfloor\theta_\rho L\rfloor$ for any fixed sufficiently small $\theta_\rho>0$ gives the simpler uniform bound
\begin{equation}\label{eq:FR-simple}
 T\geq c_\rho\frac L\delta\exp(-C_\rho D/L).
\end{equation}
Together with Proposition~\ref{prop:integrated-sign}, this independently proves Theorem~\ref{thm:joint-matched}. It also covers $L/D\to\infty$, with no growth restriction. If $K=0$, every $r$ satisfies \eqref{eq:FR-halfthreshold}; letting $r\to\infty$ in \eqref{eq:FR-all-bound} rules out, for sufficiently small $\delta$, any finite-query converter that is exact on the whole interval $[\delta,\rho\delta]$.

\subsection{A pinned construction with a growing threshold index}
We now let the order $r$ of the pinned construction grow as $\delta\downarrow0$, tracking all constants, including those raised to powers that grow with the degree.
\begin{theorem}[Uniform growing-index upper bound]\label{thm:joint-pinned}
Fix $\rho>1$ and $0<c\leq1$. Suppose $K=K(\delta)\downarrow0$ and
\[
 r=\min\{j\geq1:G_j\leq K\}\longrightarrow\infty,
 \qquad r=o(D).
\]
Then, for all sufficiently small $\delta$ along this parameter choice,
\begin{equation}\label{eq:joint-pinned-query}
 Q(\delta,K\delta;\rho,c)
 \leq C_{\rho,c}\left(r^3\delta^{-1+1/(2r)}+D\right).
\end{equation}
The constant does not depend on $r$. The construction has error at most $G_r\delta$, with no $o(\delta)$ excess, on the low band and at most $K\delta$ on the high band.
\end{theorem}
\begin{proof}
Put $n=2r$, $p=r+2$, $s=r+1$, and $P=P_r^*$ from \eqref{eq:extremizer}. The two-sided bounds
\begin{equation}\label{eq:pinned-threshold-size}
 c_\rho r^{-1}R_0^{-n}\leq G_r\leq C_\rho r^{-1}R_0^{-n}
\end{equation}
hold for every $r\geq1$ by Proposition~\ref{prop:threshold-asymptotics} and adjustment of finitely many constants. The positive contact points of $P$ are
\[
 y_j=m_\rho+h_\rho\cos(2\pi j/n),\qquad 0\leq j\leq r.
\]
Let $k$ be the least odd integer at least $A\delta^{-1+1/n}$, where the fixed $A=A_\rho$ will be chosen below. With $a_j=\arcsin(\delta y_j)$ and $\theta=\arcsin x$, define
\begin{equation}\label{eq:uniform-pinned-gate}
 W(x)=1-\prod_{j=0}^r\prod_{\sigma=\pm1}
 \left(1-D_{k,\sigma a_j}(\theta)^{2p}\right)^s,
 \qquad
 D_{k,b}(\theta)=\frac{\sin(k(\theta-b))}{k\sin(\theta-b)}.
\end{equation}
The same argument as Lemma~\ref{lem:pinned} proves that $W$ is an even polynomial, $0\leq W\leq1$, and $W(\delta y_j)=1$. More explicitly, its trigonometric degree and therefore its degree in $x$ are at most
$4(r+1)^2(r+2)(k-1)$.
Note that
\begin{equation}\label{eq:k-rounding}
 k\delta\leq2A\delta^{1/n}=o(1),\qquad
 k\geq A\delta^{-1+1/n},
\end{equation}
for sufficiently small $\delta$. These inequalities suffice; no degree-dependent constant is hidden in integer rounding.

On the low band, write $y=m_\rho+h_\rho\cos\phi$, $0\leq\phi\leq\pi$, and select a nearest angle $\phi_j=2\pi j/n$. Then $\Delta=|\phi-\phi_j|\leq\pi/n$, and $d=|y-y_j|\leq h_\rho\Delta\leq C_\rho/r$. The normalized Dirichlet kernel is the average of the exponentials with frequencies $-(k-1),-(k-3),\ldots,k-1$. Therefore $1-D_{k,0}(u)\leq k^2u^2/2$ and, since $|D|\leq1$,
$1-D_{k,0}(u)^{2p}\leq p(1-D_{k,0}(u)^2)\leq p k^2u^2$.
Keeping the factor for $y_j$ with $\sigma=+1$ and bounding the other factors of $1-W$ by one gives
\begin{equation}\label{eq:uniform-low-leakage}
 \lambda(y):=1-W(\delta y)
 \leq[C_\rho r(k\delta)^2d^2]^{r+1}.
\end{equation}
On the other hand, with $e(y)=P(y)-y$,
\begin{equation}\label{eq:uniform-contact-slack}
 G_r-e(y)=G_r[1-\cos(n\Delta)]
 \geq\frac{2G_rn^2\Delta^2}{\pi^2}
 \geq c_\rho G_rr^2d^2.
\end{equation}
Dividing \eqref{eq:uniform-low-leakage} by $\delta$ times this slack, using \eqref{eq:pinned-threshold-size}, \eqref{eq:k-rounding}, and $d\leq C_\rho/r$, yields
\begin{equation}\label{eq:uniform-slack-ratio}
 \frac{\lambda(y)}{\delta(G_r-e(y))}
 \leq C_{\rho,A}e^{-D/r}
          \left(\frac{C_{\rho,A}R_0^2}{r}\right)^r=o(1).
\end{equation}
At contacts, both numerator and denominator vanish; the ratio is defined by continuity, and the estimate remains valid. In particular $W=1-o(1)$ uniformly on the low band, and eventually
$\lambda\leq\tfrac12W\delta(G_r-e)$.

Choose $S_N$ as in \eqref{eq:SN}, and set
\begin{equation}\label{eq:joint-blend}
 q(x)=W(x)[1-\delta P(x/\delta)]+(1-W(x))S_N(x^2).
\end{equation}
For $x=\delta y$ its signed error is
$(1-x)-q(x)=W\delta e(y)+\lambda d_N(x)$, where
$d_N(x)=1-x-S_N(x^2)\in[0,1]$. The lower error bound is $-G_r\delta$. The upper bound follows from
\[
 G_r\delta-[(1-x)-q(x)]
 =W\delta(G_r-e)+\lambda(G_r\delta-d_N)\geq0.
\]
This proves that the low-band error is at most $G_r\delta$.

We next verify global contractivity, with explicit tail constants. For $|x|\geq4\rho\delta$ and small $\delta$,
\[
 |\sin(\arcsin x\pm a_j)|
 \geq |x|\sqrt{1-\rho^2\delta^2}-\rho\delta\geq |x|/2.
\]
The elementary product bound therefore gives
\begin{equation}\label{eq:uniform-tail}
 W(x)\leq\min\left\{1,\;2(r+1)^2
                \left(\frac2{k|x|}\right)^{n+4}\right\}.
\end{equation}
The factor $2^{n+4}$ cannot be absorbed into a constant independent of $r$.
For a fixed $B_\rho\geq1$, the elementary Chebyshev estimate
$|T_n(v)|\leq[2(1+|v|)]^n$ gives
\begin{equation}\label{eq:uniform-P-growth}
 0\leq P(y)\leq |y|+G_r[B_\rho(1+|y|)]^n.
\end{equation}
On $|x|\leq4\rho\delta$, this bounds $\delta P(x/\delta)$ by
$4\rho\delta+C_\rho\delta\Lambda_\rho^n/r=o(1)$, since $n=o(D)$.
On $4\rho\delta\leq|x|\leq1/k$, use $W\leq1$ and enlarge a fixed $B_\rho$ to replace $1+|x|/\delta$ by a constant times $|x|/\delta$. Then
\[
 \delta P(x/\delta)W(x)
 \leq k^{-1}+\frac{C_\rho}{r}
                     \left(\frac{B_\rho}{AR_0}\right)^n.
\]
Here $\delta(k\delta)^{-n}\leq A^{-n}$ by the lower bound in \eqref{eq:k-rounding}. Finally, on $|x|\geq1/k$, which exceeds $4\rho\delta$ eventually, combine \eqref{eq:uniform-tail} and \eqref{eq:uniform-P-growth}. The linear term is at most $C r^2 2^n/k=o(1)$ because its logarithm is at most $-D+D/n+O_\rho(n+\log r)$. The degree-$n$ term decreases as $|x|^{-4}$, and at $|x|=1/k$ is at most
\[
 C_\rho r\left(\frac{2B_\rho}{AR_0}\right)^n.
\]
Choose $A$ large enough that $2B_\rho/(AR_0)<1/2$. All three regions then give $\sup_{[-1,1]}\delta P(x/\delta)W(x)=o(1)$, in particular at most one eventually. In \eqref{eq:joint-blend} write $q=B-\delta PW$, with $B=W+(1-W)S_N\in[0,1]$. Since $P\geq0$, this proves $-1\leq q\leq1$ everywhere.

For the high band $[c,1]$, \eqref{eq:uniform-tail} implies
$W\leq C r^2(2/(kc))^{n+4}$. Separating the linear and degree-$n$ terms of \eqref{eq:uniform-P-growth} gives
\begin{align*}
 \sup_{[c,1]}\frac{W(1+|x|)}{G_r\delta}
 &\leq C_{\rho,c}r^3
       \left(\frac{C_{\rho,c}R_0}{A}\right)^n
          \delta^{n+2-4/n}=o(1),\\
 \sup_{[c,1]}\frac{W G_r\delta
                    [B_\rho(1+|x|/\delta)]^n}{G_r\delta}
 &\leq C_{\rho,c}r^2
       \left(\frac{C_{\rho,c}}{A}\right)^n
          \delta^{3-4/n}=o(1).
\end{align*}
Fixed factors $A^{-4}$ have been bounded by one. Even if a displayed base exceeds one, its logarithm is $O_{\rho,c}(n)=o(D)$, whereas the last power of $\delta$ has logarithm $-(3-o(1))D$. Thus the gate contribution $W|1-\delta P-S_N|$ is $o(G_r\delta)$, including the term $W$ itself. Choose $N$ so that the tail in \eqref{eq:SN} is at most $K\delta/2$ on the high band: for $c<1$ it suffices to take
\[
 N+1\geq\frac{\log(2/(K\delta))}{-\log(1-c^2)},
\]
and for $c=1$ take $N=1$. Since $G_r\leq K$, \eqref{eq:pinned-threshold-size} gives $L\leq na+\log r+O_\rho(1)=O_\rho(r+\log r)=o(D)$. Hence $N=O_c(D)$. The high error is at most $K\delta$ eventually, and the degree of \eqref{eq:joint-blend} is $O_\rho(r^3k+N)$. Lemma~\ref{lem:implementation} completes the proof.
\end{proof}

\subsection{What the bounds determine at shrinking accuracy}
\begin{corollary}[Bounds at intermediate accuracy]\label{cor:joint-near-match}
Under the assumptions of Theorem~\ref{thm:joint-pinned}, for all sufficiently small $\delta$,
\begin{equation}\label{eq:joint-margin-comparison}
 c_\rho r(G_{r-1}-K)^{1/r}\delta^{-1+1/(2r)}
 \leq Q(\delta,K\delta;\rho,c)
 \leq C_{\rho,c}r^3\delta^{-1+1/(2r)}.
\end{equation}
For each fixed $\gamma\in(0,1)$, if also $K\leq(1-\gamma)G_{r-1}$, then
\begin{equation}\label{eq:joint-away-boundary}
 c_{\rho,\gamma}r\delta^{-1+1/(2r)}
 \leq Q(\delta,K\delta;\rho,c)
 \leq C_{\rho,c}r^3\delta^{-1+1/(2r)}.
\end{equation}
In particular, \eqref{eq:joint-away-boundary} holds at the thresholds $K=G_r$ for all large $r$, with a fixed $\gamma$ that depends only on $\rho$.
\end{corollary}
\begin{proof}
Since $r=o(D)$, eventually $r-1\leq c_\rho\delta^{-1/2}$. Minimality of $r$ gives $K<G_{r-1}$, so Theorem~\ref{thm:FR-margin} at index $r-1$ proves the lower half of \eqref{eq:joint-margin-comparison}. Theorem~\ref{thm:joint-pinned} proves the upper half; its additive $D$ is absorbed since $r\geq2$ eventually and $D=o(\delta^{-3/4})$. If also $K\leq(1-\gamma)G_{r-1}$, then \eqref{eq:pinned-threshold-size} implies
\[
 (G_{r-1}-K)^{1/r}\geq
 \left(\frac{c_\rho\gamma}{r-1}R_0^{-2r+2}\right)^{1/r}
 \geq c_{\rho,\gamma}>0.
\]
Finally $G_r/G_{r-1}\to R_0^{-2}<1$, so any fixed $\gamma<(1-R_0^{-2})$ applies at $K=G_r$ eventually.
\end{proof}
For $K\to0$ with $L=o(D)$, Proposition~\ref{prop:threshold-asymptotics} gives $r=\Theta_\rho(L)=o(D)$. Equations \eqref{eq:FR-simple} and \eqref{eq:joint-pinned-query} then imply
\[
 \frac{\log Q(\delta,K\delta;\rho,c)}{\log(1/\delta)}\longrightarrow1.
\]
Thus the query exponent tends to one in this regime. When $K$ stays away from $G_{r-1}$, the upper end of its tier $G_r\leq K<G_{r-1}$, \eqref{eq:joint-away-boundary} bounds the ratio of the upper and lower bounds by $O_{\rho,c,\gamma}(r^2)=O_{\rho,c,\gamma}(L^2)$. Near $G_{r-1}$, our lower bound keeps the factor $(G_{r-1}-K)^{1/r}$ of \eqref{eq:joint-margin-comparison}; the independent bound \eqref{eq:FR-max} also applies there. Determining the query count up to a constant factor uniformly over this regime, including as $K$ approaches the upper end of a tier, remains open. In contrast, Theorem~\ref{thm:joint-matched} determines the order of the query count whenever $L$ is at least a fixed positive multiple of $D$.
